\documentclass{article}
\usepackage[UTF8]{ctex}
\usepackage{amsmath,amssymb,amsfonts}
\usepackage{enumitem}
\usepackage{geometry}
\geometry{a4paper, margin=1in}
% === 错题本排版优化（不参与题目标识） ===
\setlength{\parindent}{0pt}
\setlength{\parskip}{0.35em plus 0.10em minus 0.08em}
\setlength{\emergencystretch}{3em}
\setlist{itemsep=0.25em, topsep=0.40em, parsep=0pt}
\allowdisplaybreaks[2]
% === 错题本排版优化结束 ===
\begin{document}
\section*{四川大学2026年高等代数真题}
\section*{157．四川大学2026年高等代数试卷}
一．解答如下问题：\\

\subsection{题目1}
1．设 \(f(x)=x^{n}+2025 x+3\) ，其中 \(n\) 是大于 1 的整数．证明：\(f(x)\) 的全部复根互不相同．

\subsection{题目2}
2．设 \(g(x)=x^{5}-2 x^{4}+2 x^{3}-x^{2}-x+1, h(x)=x^{2026}-7\) 。证明：对于数域 \(\mathbb{F}\) 上的任意多项式 \(q(x)\)，都存在 \(\mathbb{F}\) 上的多项式 \(u(x), v(x)\) ，使得 \(q(x)=u(x) g(x)+v(x) h(x)\) 。

\subsection{题目3}
3．是否存在特征多项式为 \(x^{6}-4 x^{5}+6 x^{4}-8 x^{3}+9 x^{2}-4 x+4\) 的实对称矩阵？说明理由．

\subsection{题目4}
4．证明：多项式 \(x^{2026}+x^{2024}+x^{3}-3 x+1\) 的根不可能都是实根．
\end{document}
